\documentclass[11pt,a4paper]{article}

\usepackage[margin=1in]{geometry}
\usepackage[T1]{fontenc}
\usepackage[utf8]{inputenc}
\usepackage{lmodern}
\usepackage{microtype}
\usepackage{hyperref}
\usepackage{booktabs}
\usepackage{longtable}
\usepackage{array}
\usepackage{enumitem}
\usepackage{xcolor}
\usepackage{listings}
\usepackage{fancyhdr}
\usepackage{titlesec}
\usepackage{parskip}

\hypersetup{
  colorlinks=true,
  linkcolor=blue,
  urlcolor=blue,
  pdftitle={DTU HPC Runbook for AI Coding Agents},
  pdfauthor={OpenAI}
}

\pagestyle{fancy}
\fancyhf{}
\lhead{DTU HPC Runbook}
\rhead{\thepage}

\titleformat{\section}{\large\bfseries}{\thesection}{0.6em}{}
\titleformat{\subsection}{\normalsize\bfseries}{\thesubsection}{0.6em}{}
\titleformat{\subsubsection}{\normalsize\bfseries}{\thesubsubsection}{0.6em}{}

\lstdefinestyle{shell}{
  basicstyle=\ttfamily\small,
  breaklines=true,
  columns=fullflexible,
  frame=single,
  keepspaces=true,
  showstringspaces=false
}

\title{\textbf{DTU HPC Runbook for AI Coding Agents}}
\author{}
\date{\today}

\begin{document}
\maketitle
\tableofcontents
\newpage

\section{Purpose and scope}

This document is an operator-style reference for an AI coding agent working on DTU HPC. It is biased toward practical software engineering and ML/DL workflows rather than classical HPC simulation workflows.

It covers:

\begin{itemize}[leftmargin=1.5em]
  \item access: VPN, SSH, ThinLinc, and remote-work caveats,
  \item shell basics and working conventions,
  \item Git and GitHub authentication on the cluster,
  \item DTU software environment: modules, Python, virtual environments, \texttt{uv}, and DCC software stacks,
  \item LSF job scripts, queues, GPUs, monitoring, and failure diagnosis,
  \item storage, quota, data transfer, and scratch usage,
  \item common failure modes seen in coding/ML workflows,
  \item practical best practices for automated agents.
\end{itemize}

\section{Mental model of DTU HPC}

DTU HPC is not ``a remote Linux box where you just run things''. It is a multi-user, scheduler-managed HPC environment.

The practical model is:

\begin{enumerate}[leftmargin=1.5em]
  \item connect to DTU network (on campus or via VPN),
  \item SSH into a login node,
  \item do light editing, Git, environment setup, and inspection,
  \item use \texttt{linuxsh} for interactive preparation/tests,
  \item submit real work as LSF batch jobs with \texttt{bsub},
  \item monitor jobs and iterate.
\end{enumerate}

\textbf{Critical rule:} do not run real compute on login nodes.

\section{Access, identity, and remote-work modes}

\subsection{Access modes}

DTU documents two primary access modes:

\begin{itemize}[leftmargin=1.5em]
  \item \textbf{SSH} for command-line work.
  \item \textbf{ThinLinc} for graphical remote desktop access.
\end{itemize}

Useful endpoints:

\begin{itemize}[leftmargin=1.5em]
  \item \texttt{login1.hpc.dtu.dk}
  \item \texttt{login2.hpc.dtu.dk}
  \item \texttt{login1.gbar.dtu.dk}
  \item \texttt{login2.gbar.dtu.dk}
  \item \texttt{thinlinc.gbar.dtu.dk}
  \item \texttt{transfer.gbar.dtu.dk}
  \item VPN server: \texttt{vpn.dtu.dk}
\end{itemize}

\subsection{VPN (Cisco Secure Client / AnyConnect)}

From outside DTU, most services require VPN.

\textbf{Practical workflow:}
\begin{enumerate}[leftmargin=1.5em]
  \item install Cisco Secure Client / AnyConnect,
  \item connect to \texttt{vpn.dtu.dk},
  \item authenticate with DTU username/password,
  \item complete MFA.
\end{enumerate}

If you are already on VPN, standard password-based SSH login works. If you are outside DTU \emph{without} VPN, DTU requires SSH-key based access with an SSH key, a key passphrase, and your DTU password.

\subsection{SSH login}

Typical command:

\begin{lstlisting}[style=shell]
ssh <userid>@login1.hpc.dtu.dk
\end{lstlisting}

Alternative documented hosts:

\begin{lstlisting}[style=shell]
ssh <userid>@login2.hpc.dtu.dk
ssh <userid>@login1.gbar.dtu.dk
ssh <userid>@login2.gbar.dtu.dk
\end{lstlisting}

\subsection{Interactive shell after login}

After SSH login, move off the login node:

\begin{lstlisting}[style=shell]
linuxsh
\end{lstlisting}

That gives an interactive shell on a shared interactive node where you can:

\begin{itemize}[leftmargin=1.5em]
  \item inspect files,
  \item prepare data,
  \item compile,
  \item create environments,
  \item run small tests,
  \item submit jobs.
\end{itemize}

\textbf{Do not assume the interactive node is private.}

\subsection{ThinLinc}

ThinLinc is useful if you need GUI applications, file browsing, or GUI-heavy workflows.

On campus or on VPN:
\begin{itemize}[leftmargin=1.5em]
  \item Server: \texttt{thinlinc.gbar.dtu.dk}
  \item Username: DTU userid
  \item Password: DTU password
\end{itemize}

Outside DTU without VPN, ThinLinc must be configured to use public-key authentication, and the key must have a passphrase.

\subsection{SSH keys for access without VPN}

If you need direct external access without VPN, generate a local SSH key and follow DTU's outside-access guide.

Practical command on your own machine:

\begin{lstlisting}[style=shell]
ssh-keygen -t ed25519 -f ~/.ssh/gbar
\end{lstlisting}

\textbf{Always verify host fingerprints on first connection.}

\subsection{Recommended local SSH config}

On your local machine:

\begin{lstlisting}[style=shell]
Host dtu-hpc
  HostName login1.hpc.dtu.dk
  User <userid>
  ServerAliveInterval 30
  ServerAliveCountMax 4
\end{lstlisting}

Then connect with:

\begin{lstlisting}[style=shell]
ssh dtu-hpc
\end{lstlisting}

\section{Working conventions for an AI coding agent}

\subsection{Core rules}

\begin{enumerate}[leftmargin=1.5em]
  \item Never do heavy compute on login nodes.
  \item Prefer stable absolute paths for outputs, checkpoints, and databases.
  \item Install environments once; do not rebuild environments inside every job.
  \item Submit from repo root unless you know exactly what the job CWD will be.
  \item Treat quota and filesystem pressure as first-class failure sources.
  \item Use short smoke tests before expensive runs.
  \item Log aggressively and unbuffer stdout for debugging.
\end{enumerate}

\subsection{Suggested repo layout}

A practical project layout:

\begin{lstlisting}[style=shell]
repo/
  .venv/
  logs/
  submit_scripts/
  src/
  data/
  runs/
  scripts/
  pyproject.toml
  uv.lock
\end{lstlisting}

\subsection{Recommended shell habits}

\begin{lstlisting}[style=shell]
pwd
ls -la
whoami
hostname
git status
module list
python --version
which python
\end{lstlisting}

\section{Filesystems, quota, and data transfer}

\subsection{Home directory}

Each user gets a default \texttt{zhome} quota of 30 GB. The home directory is backed up.

Check usage:

\begin{lstlisting}[style=shell]
getquota_zhome.sh
\end{lstlisting}

\subsection{Scratch / work space}

If you need temporary large space, ask DTU support for scratch. DTU documents scratch under \texttt{/work3/\$USER}. Scratch is not backed up.

Check scratch usage:

\begin{lstlisting}[style=shell]
getquota_work3.sh
\end{lstlisting}

\subsection{What to store where}

\begin{longtable}{p{0.27\linewidth} p{0.67\linewidth}}
\toprule
Location & Recommended contents \\
\midrule
\texttt{zhome} & code, configs, small outputs, important scripts, notebooks, lightweight metadata, anything you need backed up \\
\texttt{/work3} or scratch & large datasets, caches, checkpoints, temporary feature banks, Optuna databases, model artifacts, Hugging Face cache, W\&B cache \\
repo-local \texttt{runs/} & fine for small/medium outputs; not ideal for huge artifacts if quota is tight \\
\bottomrule
\end{longtable}

\subsection{Large data transfer}

For large transfers, DTU recommends using:

\begin{lstlisting}[style=shell]
transfer.gbar.dtu.dk
\end{lstlisting}

Examples:

\begin{lstlisting}[style=shell]
scp -i ~/.ssh/gbar file1.txt s123456@transfer.gbar.dtu.dk:
scp -i ~/.ssh/gbar s123456@transfer.gbar.dtu.dk:file1.txt .
rsync -avP localdir/ s123456@transfer.gbar.dtu.dk:remote-dir/
\end{lstlisting}

\subsection{Cleaning quota pressure}

Common space hogs on DTU HPC:

\begin{itemize}[leftmargin=1.5em]
  \item \texttt{\$HOME/.local/share/Trash}
  \item \texttt{\$HOME/.cache/pip}
  \item \texttt{\$HOME/.cache/huggingface}
  \item \texttt{\$HOME/.cache/wandb}
  \item \texttt{\$HOME/.vscode-server}
  \item local package caches for pip/conda/uv
\end{itemize}

Useful commands:

\begin{lstlisting}[style=shell]
du -h --apparent --max-depth=1 .
pip cache purge
uv cache clean
rm -rvf $HOME/.cache/huggingface/*
rm -rvf $HOME/.local/share/Trash/*
\end{lstlisting}

\section{Shell and Bash basics an agent should know}

\subsection{Navigation and file inspection}

\begin{lstlisting}[style=shell]
pwd
ls -la
cd /path/to/dir
mkdir -p logs
cp src dst
mv src dst
rm -rf path
find . -name "*.py"
du -h --apparent --max-depth=1 .
df -h .
\end{lstlisting}

\subsection{Permissions and executability}

\begin{lstlisting}[style=shell]
chmod +x script.sh
ls -l script.sh
\end{lstlisting}

\subsection{Line endings}

If a script was edited on Windows and fails on Linux, normalize line endings:

\begin{lstlisting}[style=shell]
sed -i 's/\r$//' script.sh
\end{lstlisting}

\subsection{Environment variables}

\begin{lstlisting}[style=shell]
echo $PATH
export PYTHONUNBUFFERED=1
export HF_HOME=/work3/$USER/huggingface
export WANDB_DIR=/work3/$USER/wandb
\end{lstlisting}

\section{Git and GitHub on DTU HPC}

\subsection{Clone a repo}

Create a workspace:

\begin{lstlisting}[style=shell]
mkdir -p ~/projects
cd ~/projects
git clone <REPO_URL>
cd <repo-name>
\end{lstlisting}

\subsection{Git branch operations}

\begin{lstlisting}[style=shell]
git branch -a
git switch <branch>
git fetch
git switch -c <branch> --track origin/<branch>
git pull
\end{lstlisting}

\subsection{Force-sync to remote and discard local changes}

\begin{lstlisting}[style=shell]
git fetch origin
git reset --hard origin/<branch>
git clean -fd      # remove untracked files
git clean -fdx     # also remove ignored files
\end{lstlisting}

\subsection{GitHub authentication: important}

GitHub no longer supports account-password authentication for Git over HTTPS. You must use:

\begin{itemize}[leftmargin=1.5em]
  \item a Personal Access Token (PAT) over HTTPS, or
  \item SSH keys.
\end{itemize}

\subsection{HTTPS with PAT}

\begin{lstlisting}[style=shell]
export GIT_ASKPASS=
export SSH_ASKPASS=
export DISPLAY=
git clone https://github.com/<user>/<repo>.git
\end{lstlisting}

When prompted:
\begin{itemize}[leftmargin=1.5em]
  \item username = GitHub username
  \item password = PAT
\end{itemize}

\subsection{SSH to GitHub (recommended)}

Generate a key on the cluster:

\begin{lstlisting}[style=shell]
ssh-keygen -t ed25519 -a 64 -f ~/.ssh/id_ed25519
cat ~/.ssh/id_ed25519.pub
\end{lstlisting}

Add the public key to GitHub, then:

\begin{lstlisting}[style=shell]
git clone git@github.com:<user>/<repo>.git
\end{lstlisting}

\subsection{Existing SSH keys}

\begin{lstlisting}[style=shell]
ls -al ~/.ssh
\end{lstlisting}

\section{VS Code Remote-SSH}

DTU explicitly says VS Code Remote is unsupported, though many users use it anyway.

\subsection{What works in practice}

\begin{itemize}[leftmargin=1.5em]
  \item Remote-SSH into \texttt{login1.hpc.dtu.dk} or \texttt{login2.hpc.dtu.dk}.
  \item Use VS Code for editing, Git, file browsing, and terminal workflows.
  \item Use the terminal to run \texttt{linuxsh} and to submit jobs.
\end{itemize}

\subsection{What to avoid}

\begin{itemize}[leftmargin=1.5em]
  \item Do not treat a login-node VS Code terminal as a compute node.
  \item Do not leave large training or preprocessing jobs running directly in a VS Code terminal on a login node.
\end{itemize}

\subsection{Known incompatibility}

If VS Code Remote breaks with:

\begin{lstlisting}[style=shell]
The remote host may not meet VS Code Server's prerequisites
for glibc and libstdc++
\end{lstlisting}

DTU's note says to downgrade the client to version 1.85.

\section{Modules and software stacks}

\subsection{Why modules exist}

The cluster provides many software versions. Modules let you choose the versions you want without polluting the global environment.

\subsection{Core commands}

\begin{lstlisting}[style=shell]
module list
module avail
module load <name>
module unload <name>
module switch <old> <new>
module whatis <name>
module show <name>
\end{lstlisting}

\subsection{Important behavior}

\begin{itemize}[leftmargin=1.5em]
  \item Your login environment is intentionally minimal.
  \item Most scientific software is \emph{not} available until you load modules.
  \item Batch jobs should load all needed modules explicitly.
\end{itemize}

\subsection{DCC software stack}

DTU also provides a DCC software stack. Enable it with:

\begin{lstlisting}[style=shell]
source /dtu/sw/dcc/dcc-sw.bash
\end{lstlisting}

Important caveats:

\begin{itemize}[leftmargin=1.5em]
  \item enabling the DCC stack disables the regular modules,
  \item for reproducibility, use a fixed \texttt{dcc-setup/<version>} instead of implicitly taking ``latest''.
\end{itemize}

Example:

\begin{lstlisting}[style=shell]
module purge
module load dcc-setup/2023-aug
\end{lstlisting}

\section{Python, virtual environments, and \texttt{uv}}

\subsection{Python on DTU HPC}

DTU's Python guidance boils down to:

\begin{itemize}[leftmargin=1.5em]
  \item choose a Python version via modules,
  \item do not mix packages built for different Python versions,
  \item prefer virtual environments,
  \item use \texttt{python3 -m pip ...} rather than raw \texttt{pip3 install ...},
  \item do not mix conda and pip unless you really know what you are doing.
\end{itemize}

\subsection{Find and load Python}

\begin{lstlisting}[style=shell]
module avail python
module load python3/3.10.13
python3 --version
which python3
\end{lstlisting}

\subsection{Classic venv workflow}

\begin{lstlisting}[style=shell]
module load python3/3.10.13
python3 -m venv .venv
. .venv/bin/activate
python3 -m pip install -U pip
python3 -m pip install -r requirements.txt
\end{lstlisting}

\subsection{\texttt{uv} workflow}

If you use \texttt{uv}, do it \emph{on the cluster}. Do not reuse a Windows-created \texttt{.venv} on Linux.

Typical pattern:

\begin{lstlisting}[style=shell]
module load python3/3.10.13
python3 -m pip install --user -U uv
export PATH="$HOME/.local/bin:$PATH"

uv venv .venv --python "$(which python3)"
uv sync --frozen
\end{lstlisting}

Then in jobs, prefer:

\begin{lstlisting}[style=shell]
export PATH="$HOME/.local/bin:$PATH"
uv run --no-sync -m mypackage.entrypoint
\end{lstlisting}

\subsection{Install once, reuse everywhere}

Because the filesystem is shared across nodes, install the environment once (typically from an interactive node) and reuse it from batch jobs.

\subsection{PyTorch notes}

DTU's PyTorch page says:

\begin{itemize}[leftmargin=1.5em]
  \item PyTorch is not preinstalled.
  \item Install it in a virtual environment in home or scratch.
  \item In batch jobs, \texttt{CUDA\_VISIBLE\_DEVICES} is set automatically by the scheduler.
\end{itemize}

Useful sanity checks:

\begin{lstlisting}[style=shell]
nvidia-smi
python -c "import torch; print(torch.cuda.is_available())"
python -c "import torch; print(torch.cuda.get_device_name(0))"
\end{lstlisting}

\section{LSF scheduler: workflow, concepts, and commands}

\subsection{Core lifecycle}

\begin{enumerate}[leftmargin=1.5em]
  \item write a submit script,
  \item request resources with \texttt{\#BSUB} lines,
  \item submit with \texttt{bsub < script.sh},
  \item monitor with \texttt{bstat}, \texttt{bpeek}, \texttt{showstart}, \texttt{bhist}, \texttt{bacct},
  \item kill with \texttt{bkill} if needed.
\end{enumerate}

\subsection{Key command summary}

\begin{longtable}{p{0.24\linewidth} p{0.70\linewidth}}
\toprule
Command & Meaning \\
\midrule
\texttt{bsub < job.sh} & submit a job \\
\texttt{bstat} & DTU compact status view \\
\texttt{bstat -C} & show CPU efficiency \\
\texttt{bstat -M} & show memory usage/limit \\
\texttt{bkill <jobid>} & kill job \\
\texttt{bkill -s SIGTERM <jobid>} & request graceful shutdown \\
\texttt{classstat <queue>} & queue-level load summary \\
\texttt{nodestat <queue>} & node-level availability summary \\
\texttt{nodestat -F <queue>} & show CPU model/features/memory \\
\texttt{nodestat -G <queue>} & GPU details \\
\texttt{showstart <jobid>} & estimate start time for pending jobs \\
\texttt{bpeek <jobid>} & view partial stdout/stderr of running job \\
\texttt{bhist -l <jobid>} & detailed job history / postmortem \\
\texttt{bacct <jobid>} & finished-job accounting summary \\
\bottomrule
\end{longtable}

\subsection{Status values}

From DTU's \texttt{bstat} page:

\begin{itemize}[leftmargin=1.5em]
  \item \texttt{PEND}: queued, waiting for scheduling
  \item \texttt{RUN}: running
  \item \texttt{DONE}: completed successfully
  \item \texttt{EXIT}: exited/killed/failed
  \item \texttt{SSUSP}: suspended
\end{itemize}

In DTU's compact output, \texttt{NALLOC} is the number shown in the allocation column and in practice corresponds to the allocated slots/cores for the job.

\subsection{General job-script structure}

\begin{lstlisting}[style=shell]
#!/bin/sh
#BSUB -q hpc
#BSUB -J myjob
#BSUB -n 4
#BSUB -R "span[hosts=1]"
#BSUB -R "rusage[mem=4GB]"
#BSUB -M 5GB
#BSUB -W 24:00
#BSUB -o Output_%J.out
#BSUB -e Output_%J.err

module load <modules>
my_command_here
\end{lstlisting}

\subsection{Meaning of common \texttt{\#BSUB} directives}

\begin{longtable}{p{0.24\linewidth} p{0.70\linewidth}}
\toprule
Directive & Meaning \\
\midrule
\texttt{-q} & queue name \\
\texttt{-J} & job name \\
\texttt{-n} & total CPU slots/cores \\
\texttt{-R "span[hosts=1]"} & keep all requested CPU slots on one host \\
\texttt{-R "span[ptile=N]"} & spread slots in groups of N per host; useful for MPI-style layouts \\
\texttt{-R "rusage[mem=5GB]"} & memory request \emph{per core/slot} \\
\texttt{-M 5GB} & per-process memory limit; exceeding host-enforced limit can kill the job \\
\texttt{-W 12:00} & walltime limit \\
\texttt{-B} & mail at start \\
\texttt{-N} & mail at completion \\
\texttt{-Ne} & mail on failure only \\
\texttt{-o} & stdout file \\
\texttt{-e} & stderr file \\
\texttt{-gpu "num=1:mode=exclusive\_process"} & request GPUs in exclusive-process mode \\
\bottomrule
\end{longtable}

\subsection{CPU and memory reasoning}

Important consequence of DTU's memory model:

\begin{itemize}[leftmargin=1.5em]
  \item \texttt{rusage[mem=5GB]} with \texttt{-n 4} means about 20 GB total RAM requested.
  \item memory limits are effectively enforced per host, not globally across arbitrarily many hosts.
\end{itemize}

\subsection{Queue defaults and limits}

DTU documents for the main CPU queue family:

\begin{itemize}[leftmargin=1.5em]
  \item default queue: \texttt{hpc}
  \item max wallclock: 72 hours
  \item default walltime: 15 minutes
\end{itemize}

ThinLinc/application nodes differ and are not batch-compute targets in the same way.

\section{GPU queues and hardware}

\subsection{Current public GPU queues}

DTU's public GPU page currently lists:

\begin{longtable}{p{0.20\linewidth} p{0.58\linewidth} p{0.14\linewidth}}
\toprule
Queue & Hardware & Notes \\
\midrule
\texttt{gpul40s} & 6 nodes, each with 2x Nvidia L40S PCIe 48 GB & public GPU queue \\
\texttt{gpua100} & 4 nodes with 2x A100 40 GB and 6 nodes with 2x A100 80 GB & A100 queue \\
\texttt{gpuv100} & 6 nodes with 2x V100 16 GB; 8 nodes with 2x V100 32 GB; 3 nodes with 4x V100 32 GB NVLink & V100 queue \\
\texttt{gpua10} & 1 node with 2x Tesla A10 24 GB & small queue \\
\texttt{gpua40} & 1 node with 2x Tesla A40 48 GB NVLink & small queue \\
\texttt{gpuamd} & AMD MI25 queue & non-Nvidia \\
\bottomrule
\end{longtable}

DTU also documents interactive GPU helper commands:

\begin{lstlisting}[style=shell]
voltash
sxm2sh
a100sh
\end{lstlisting}

\subsection{GPU selectors}

Valid selectors depend on queue.

Examples:
\begin{itemize}[leftmargin=1.5em]
  \item \texttt{gpua100} + \texttt{select[gpu40gb]}
  \item \texttt{gpua100} + \texttt{select[gpu80gb]}
  \item \texttt{gpuv100} + \texttt{select[gpu32gb]}
  \item \texttt{gpuv100} + \texttt{select[sxm2]}
\end{itemize}

\textbf{Do not mismatch queue and selector.} Example: \texttt{gpuv100} with \texttt{select[gpu40gb]} is a bad request.

\subsection{A100 compatibility}

DTU notes that code targeting Nvidia A100 should be built with CUDA 11.0 or newer.

\subsection{Choosing a queue for testing}

Practical guidance:

\begin{itemize}[leftmargin=1.5em]
  \item shortest queue is dynamic; check with \texttt{classstat} and \texttt{nodestat},
  \item cheapest/worst GPU is often \texttt{gpua10}, but it has very little capacity,
  \item \texttt{gpuv100} often starts sooner due to more nodes,
  \item use short walltimes for smoke tests.
\end{itemize}

\subsection{Current CPU hardware summary}

DTU's central cluster page currently shows a mixed CPU fleet including, among others:

\begin{itemize}[leftmargin=1.5em]
  \item 2x AMD EPYC 9554 nodes with 1536 GB RAM and EDR InfiniBand,
  \item Intel Xeon Gold 6342 nodes with 512 GB RAM,
  \item Intel Xeon Gold 6226R nodes with 384 or 768 GB RAM,
  \item Intel Xeon Gold 6126 / 6142 nodes,
  \item older Xeon E5-2650v4 / E5-2660v3 nodes,
  \item dedicated AMD queues such as \texttt{epyc}, \texttt{rome}, and \texttt{milan}.
\end{itemize}

Therefore, identical CPU jobs may run at different speeds on different hosts. If node features matter, use \texttt{nodestat -F <queue>} and explicit \texttt{select[...]} requirements.

\section{Submit-script templates}

\subsection{Minimal CPU batch script}

\begin{lstlisting}[style=shell]
#!/bin/sh
#BSUB -J cpu-job
#BSUB -o logs/%J.out
#BSUB -e logs/%J.err
#BSUB -q hpc
#BSUB -n 4
#BSUB -R "span[hosts=1]"
#BSUB -R "rusage[mem=4GB]"
#BSUB -W 02:00
#BSUB -B
#BSUB -N

cd "${LS_SUBCWD:-$PWD}"
mkdir -p logs

module load python3/3.10.13
. .venv/bin/activate

export PYTHONUNBUFFERED=1
export PYTHONPATH="$PWD:$PYTHONPATH"

python -m mypackage.main "$@"
\end{lstlisting}

\subsection{Minimal single-GPU batch script}

\begin{lstlisting}[style=shell]
#!/bin/sh
#BSUB -J gpu-job
#BSUB -o logs/%J.out
#BSUB -e logs/%J.err
#BSUB -q gpua100
#BSUB -n 4
#BSUB -gpu "num=1:mode=exclusive_process"
#BSUB -R "select[gpu40gb]"
#BSUB -R "rusage[mem=5GB]"
#BSUB -R "span[hosts=1]"
#BSUB -W 12:00
#BSUB -B
#BSUB -N

cd "${LS_SUBCWD:-$PWD}"
mkdir -p logs

. .venv/bin/activate
export PYTHONUNBUFFERED=1
export PYTHONPATH="$PWD:$PYTHONPATH"

python -m mypackage.train "$@"
\end{lstlisting}

\subsection{Single-GPU with \texttt{uv}}

\begin{lstlisting}[style=shell]
#!/bin/sh
#BSUB -J gpu-uv-job
#BSUB -o logs/%J.out
#BSUB -e logs/%J.err
#BSUB -q gpuv100
#BSUB -n 4
#BSUB -gpu "num=1:mode=exclusive_process"
#BSUB -R "select[gpu32gb]"
#BSUB -R "rusage[mem=5GB]"
#BSUB -R "span[hosts=1]"
#BSUB -W 04:00

cd "${LS_SUBCWD:-$PWD}"
mkdir -p logs

export PATH="$HOME/.local/bin:$PATH"
export PYTHONUNBUFFERED=1
export PYTHONPATH="$PWD:$PYTHONPATH"

uv run --no-sync -m mypackage.entrypoint "$@"
\end{lstlisting}

\subsection{Run another shell script from the submit script}

This is valid and often cleaner than putting everything directly in the LSF script.

\begin{lstlisting}[style=shell]
#!/bin/sh
#BSUB -J pipeline
#BSUB -o logs/%J.out
#BSUB -e logs/%J.err
#BSUB -q gpua100
#BSUB -n 4
#BSUB -gpu "num=1:mode=exclusive_process"
#BSUB -R "select[gpu40gb]"
#BSUB -R "rusage[mem=5GB]"
#BSUB -R "span[hosts=1]"
#BSUB -W 24:00

cd "${LS_SUBCWD:-$PWD}"
mkdir -p logs
export PATH="$HOME/.local/bin:$PATH"

bash scripts/run_pipeline.sh "$@"
\end{lstlisting}

\subsection{PyTorch \texttt{torchrun} single-node pattern}

\begin{lstlisting}[style=shell]
#!/bin/sh
#BSUB -J torch-single
#BSUB -q gpuv100
#BSUB -n 4
#BSUB -R "span[hosts=1]"
#BSUB -gpu "num=1:mode=exclusive_process"
#BSUB -W 01:00
#BSUB -R "rusage[mem=5GB]"
#BSUB -o logs/%J.out
#BSUB -e logs/%J.err

cd "${LS_SUBCWD:-$PWD}"
mkdir -p logs
. .venv/bin/activate

torchrun --standalone --nproc_per_node=1 train.py
\end{lstlisting}

\subsection{Multi-node / multi-GPU pattern}

Use only if your code is actually distributed-aware.

\begin{lstlisting}[style=shell]
#!/bin/sh
#BSUB -J torch-multi
#BSUB -q gpuv100
#BSUB -n 8
#BSUB -R "span[ptile=4]"
#BSUB -gpu "num=2:mode=exclusive_process"
#BSUB -W 01:00
#BSUB -R "rusage[mem=5GB]"
#BSUB -o logs/%J.out
#BSUB -e logs/%J.err

cd "${LS_SUBCWD:-$PWD}"
mkdir -p logs
. .venv/bin/activate

List=$(echo "$LSB_HOSTS" | tr ' ' '\n' | sort -u | paste -sd,)
PORT=29500
blaunch -z "$List" torchrun \
  --nnodes=2 \
  --nproc_per_node=2 \
  --rdzv_id="$LSB_JOBID" \
  --rdzv_backend=c10d \
  --rdzv_endpoint="$HOSTNAME:$PORT" \
  train.py
\end{lstlisting}

\section{Submission rules that prevent stupid failures}

\subsection{Submit from repo root}

Relative paths in jobs resolve from the job's working directory, not from the location of the script file. Prefer:

\begin{lstlisting}[style=shell]
cd "${LS_SUBCWD:-$PWD}"
\end{lstlisting}

and submit while standing in repo root:

\begin{lstlisting}[style=shell]
cd ~/projects/myrepo
mkdir -p logs
bsub < submit_scripts/train.sh
\end{lstlisting}

\subsection{Create log directories before submission}

Practical note: if your output paths are \texttt{logs/\%J.out} and \texttt{logs/\%J.err}, create \texttt{logs/} before \texttt{bsub}. LSF may open output files before your shell script body runs.

\begin{lstlisting}[style=shell]
mkdir -p logs
bsub < submit.sh
\end{lstlisting}

\subsection{Windows venvs do not work on Linux}

Never copy a Windows-created \texttt{.venv} and expect it to work on DTU HPC. Build the environment on the cluster.

\subsection{Namespace packages and \texttt{src.*}}

If your repo uses a \texttt{src/} layout and you run:

\begin{lstlisting}[style=shell]
python -m src.something.entrypoint
\end{lstlisting}

make sure the repo root is on \texttt{PYTHONPATH}:

\begin{lstlisting}[style=shell]
export PYTHONPATH="$PWD:$PYTHONPATH"
\end{lstlisting}

Alternative: install the package in editable mode:

\begin{lstlisting}[style=shell]
uv pip install -e .
# or
python -m pip install -e .
\end{lstlisting}

\section{Monitoring, diagnosis, and error interpretation}

\subsection{Check queue state before submitting}

\begin{lstlisting}[style=shell]
classstat hpc
classstat gpuv100
classstat gpua100
\end{lstlisting}

\subsection{Check node features and queue fit}

\begin{lstlisting}[style=shell]
nodestat hpc
nodestat -F hpc
nodestat -G gpua100
\end{lstlisting}

\subsection{See job state}

\begin{lstlisting}[style=shell]
bstat
bstat <jobid>
\end{lstlisting}

\subsection{Estimate start time}

\begin{lstlisting}[style=shell]
showstart <jobid>
\end{lstlisting}

\subsection{Check efficiency and memory}

\begin{lstlisting}[style=shell]
bstat -C <jobid>
bstat -M <jobid>
\end{lstlisting}

\subsection{Peek at live output}

\begin{lstlisting}[style=shell]
bpeek <jobid>
\end{lstlisting}

\subsection{Postmortem on finished/failed jobs}

\begin{lstlisting}[style=shell]
bhist -l <jobid>
bacct <jobid>
\end{lstlisting}

\subsection{Stop jobs}

\begin{lstlisting}[style=shell]
bkill <jobid>
bkill -s SIGTERM <jobid>
bkill -s SIGKILL <jobid>
\end{lstlisting}

\subsection{Useful failure patterns}

\subsubsection{Job stuck in \texttt{PEND}}

Most common causes:
\begin{itemize}[leftmargin=1.5em]
  \item queue is full,
  \item selector and queue mismatch,
  \item over-requested walltime,
  \item over-requested memory,
  \item over-requested GPUs,
  \item special queue access restriction.
\end{itemize}

Check:
\begin{lstlisting}[style=shell]
classstat <queue>
showstart <jobid>
nodestat -F <queue>
\end{lstlisting}

\subsubsection{Job goes to \texttt{EXIT}}

Most common causes:
\begin{itemize}[leftmargin=1.5em]
  \item exceeded memory limit,
  \item bad module path or missing environment,
  \item import error due to wrong working directory / \texttt{PYTHONPATH},
  \item bad output path,
  \item disk/quota failure,
  \item crashed application.
\end{itemize}

Check:
\begin{lstlisting}[style=shell]
bstat -M <jobid>
bhist -l <jobid>
bacct <jobid>
\end{lstlisting}

\subsubsection{\texttt{No module named ...}}

Usual causes:
\begin{itemize}[leftmargin=1.5em]
  \item wrong working directory,
  \item repo root not on \texttt{PYTHONPATH},
  \item package not installed editable,
  \item using a different Python than the one used to build the environment.
\end{itemize}

Minimal safe fix:
\begin{lstlisting}[style=shell]
cd "${LS_SUBCWD:-$PWD}"
export PYTHONPATH="$PWD:$PYTHONPATH"
\end{lstlisting}

\subsubsection{SQLite / Optuna \texttt{disk I/O error}}

This is common on HPC if:
\begin{itemize}[leftmargin=1.5em]
  \item your home quota is full,
  \item the filesystem is under trouble,
  \item the DB is on a bad path,
  \item multiple jobs are writing to the same SQLite DB at once.
\end{itemize}

What to do:
\begin{lstlisting}[style=shell]
getquota_zhome.sh
df -h .
find . -name optuna.db
sqlite3 /path/to/optuna.db "select max(number), count(*) from trials;"
\end{lstlisting}

Best practice:
\begin{itemize}[leftmargin=1.5em]
  \item use an absolute path for Optuna storage,
  \item put heavy DB/checkpoint artifacts on scratch if quota is tight,
  \item avoid concurrent writers to the same SQLite file,
  \item enable regular application-level checkpointing.
\end{itemize}

\subsubsection{GitHub auth failure on clone/pull}

If you see auth failures on HTTPS, do not enter your GitHub account password. Use a PAT or SSH key.

\subsubsection{VS Code glibc/libstdc++ error}

Pin VS Code to version 1.85 or stop using VS Code Remote on DTU HPC.

\section{Practical AI / ML best practices on DTU HPC}

\subsection{Environment setup}

\begin{itemize}[leftmargin=1.5em]
  \item create the environment once on an interactive node,
  \item do not install packages in the job script,
  \item keep caches off home when they get large,
  \item avoid mixing conda and pip unless unavoidable.
\end{itemize}

\subsection{Data layout}

\begin{itemize}[leftmargin=1.5em]
  \item use a small number of large files rather than huge directories with countless tiny files,
  \item do not dump millions of artifacts into a single folder,
  \item store scratchable data outside backed-up home space when possible.
\end{itemize}

\subsection{GPU usage}

\begin{itemize}[leftmargin=1.5em]
  \item verify that the code actually uses the GPU before submitting long runs,
  \item do not request multiple GPUs unless the code is multi-GPU aware,
  \item request GPU exclusive mode,
  \item align \texttt{torchrun} / distributed process counts with requested GPUs.
\end{itemize}

\subsection{CPU threads}

Many Python libraries choose thread counts badly. Make the job explicit:

\begin{lstlisting}[style=shell]
export OMP_NUM_THREADS=4
export MKL_NUM_THREADS=4
export OPENBLAS_NUM_THREADS=4
export NUMEXPR_NUM_THREADS=4
\end{lstlisting}

Set these coherently with \texttt{\#BSUB -n}.

\subsection{Checkpointing}

Long jobs should checkpoint.
If the queue enforces a 24-hour walltime for GPU jobs, treat every job as preemptible from the application point of view.

\subsection{Absolute paths for persistent state}

For study DBs, checkpoints, and outputs that must survive job restarts, prefer stable absolute paths.

Example:
\begin{lstlisting}[style=shell]
/work3/$USER/myproject/checkpoints/
/work3/$USER/myproject/optuna/bo_final.db
\end{lstlisting}

\subsection{Unbuffered logs}

Always set:

\begin{lstlisting}[style=shell]
export PYTHONUNBUFFERED=1
\end{lstlisting}

Otherwise debugging with \texttt{bpeek} is worse.

\subsection{Smoke tests}

Before submitting a 12--24 hour job, run:

\begin{itemize}[leftmargin=1.5em]
  \item 1 epoch,
  \item 1 seed,
  \item reduced dataset,
  \item 5--10 minute walltime,
  \item smallest acceptable GPU queue.
\end{itemize}

\section{Quick-reference command appendix}

\subsection{Access}

\begin{lstlisting}[style=shell]
ssh <userid>@login1.hpc.dtu.dk
linuxsh
exit
\end{lstlisting}

\subsection{Modules}

\begin{lstlisting}[style=shell]
module avail
module list
module load python3/3.10.13
module purge
\end{lstlisting}

\subsection{Environment}

\begin{lstlisting}[style=shell]
python3 -m venv .venv
. .venv/bin/activate
python3 -m pip install -U pip
uv sync --frozen
\end{lstlisting}

\subsection{Git}

\begin{lstlisting}[style=shell]
git clone <url>
git switch <branch>
git fetch
git reset --hard origin/<branch>
git clean -fdx
\end{lstlisting}

\subsection{Submission and monitoring}

\begin{lstlisting}[style=shell]
mkdir -p logs
bsub < submit.sh
bstat
bstat -M
bpeek <jobid>
bhist -l <jobid>
bacct <jobid>
bkill <jobid>
\end{lstlisting}

\subsection{Storage}

\begin{lstlisting}[style=shell]
getquota_zhome.sh
getquota_work3.sh
du -h --apparent --max-depth=1 .
df -h .
\end{lstlisting}

\subsection{Transfer}

\begin{lstlisting}[style=shell]
scp file s123456@transfer.gbar.dtu.dk:
rsync -avP localdir/ s123456@transfer.gbar.dtu.dk:remote-dir/
\end{lstlisting}

\section{Agent checklist before launching a long run}

\begin{enumerate}[leftmargin=1.5em]
  \item Am I on an interactive node or submitting to LSF rather than computing on login?
  \item Does the repo root contain the intended code branch and files?
  \item Does \texttt{python --version} match the environment I installed?
  \item Does the script set \texttt{cd "\$\{LS\_SUBCWD:-\$PWD\}"}?
  \item Does the script export \texttt{PYTHONPATH="\$PWD:\$PYTHONPATH"} if using \texttt{src.*} imports?
  \item Did I create \texttt{logs/} before \texttt{bsub}?
  \item Is the queue-selector combination valid?
  \item Is the walltime realistic?
  \item Is memory request realistic?
  \item Are checkpoints and DBs written to stable paths?
  \item Is quota sufficient?
  \item Did I run a short smoke test?
\end{enumerate}

\section{Support and source list}

\subsection{Support}

\begin{itemize}[leftmargin=1.5em]
  \item DTU HPC support: \texttt{support@hpc.dtu.dk}
  \item Pages also list \texttt{support@cc.dtu.dk} in several places
\end{itemize}

\subsection{Primary DTU references}

\begin{itemize}[leftmargin=1.5em]
  \item \url{https://www.hpc.dtu.dk/?page_id=2534}
  \item \url{https://www.hpc.dtu.dk/?page_id=2501}
  \item \url{https://www.hpc.dtu.dk/?page_id=1416}
  \item \url{https://www.hpc.dtu.dk/?page_id=1519}
  \item \url{https://www.hpc.dtu.dk/?page_id=282}
  \item \url{https://www.hpc.dtu.dk/?page_id=2008}
  \item \url{https://www.hpc.dtu.dk/?page_id=3098}
  \item \url{https://www.hpc.dtu.dk/?page_id=2759}
  \item \url{https://www.hpc.dtu.dk/?page_id=2520}
  \item \url{https://www.hpc.dtu.dk/?page_id=5685}
  \item \url{https://www.hpc.dtu.dk/?page_id=4317}
  \item \url{https://www.hpc.dtu.dk/?page_id=4377}
  \item \url{https://www.hpc.dtu.dk/?page_id=927}
  \item \url{https://www.hpc.dtu.dk/?page_id=5556}
  \item \url{https://www.hpc.dtu.dk/?page_id=3678}
  \item \url{https://www.hpc.dtu.dk/?page_id=3246}
  \item \url{https://www.hpc.dtu.dk/?page_id=3270}
  \item \url{https://www.hpc.dtu.dk/?page_id=4788}
  \item \url{https://www.hpc.dtu.dk/?page_id=4800}
  \item \url{https://www.hpc.dtu.dk/?page_id=4204}
  \item \url{https://www.hpc.dtu.dk/?p=5112}
\end{itemize}

\subsection{GitHub references}

\begin{itemize}[leftmargin=1.5em]
  \item \url{https://docs.github.com/en/authentication/keeping-your-account-and-data-secure/about-authentication-to-github}
  \item \url{https://docs.github.com/en/authentication/keeping-your-account-and-data-secure/managing-your-personal-access-tokens}
  \item \url{https://docs.github.com/en/authentication/connecting-to-github-with-ssh}
  \item \url{https://docs.github.com/en/authentication/connecting-to-github-with-ssh/generating-a-new-ssh-key-and-adding-it-to-the-ssh-agent}
  \item \url{https://docs.github.com/en/repositories/creating-and-managing-repositories/troubleshooting-cloning-errors}
\end{itemize}

\subsection{Final caution}

DTU's site is extensive but not perfectly uniform; some pages are older than others, and hardware/software details evolve. When behavior matters more than documentation, trust in this order:

\begin{enumerate}[leftmargin=1.5em]
  \item current queue/node state from \texttt{classstat} / \texttt{nodestat},
  \item current module availability from \texttt{module avail},
  \item current filesystem/quota state from \texttt{getquota\_*.sh} and \texttt{df -h},
  \item current job behavior from \texttt{bstat}, \texttt{bhist}, and \texttt{bacct},
  \item then the static documentation.
\end{enumerate}

\end{document}